\documentclass{ctexart}
\usepackage{amsmath, amssymb}
\usepackage{geometry}
\geometry{a4paper, margin=2cm}

\title{\textbf{第2章 牛顿运动定律} --- 知识点与公式}
\author{}
\date{}

\begin{document}
\maketitle

\section*{一、牛顿运动三定律}

\subsection*{1. 牛顿第一定律（惯性定律）}
任何质点都保持静止或匀速直线运动状态，直到其它物体作用的力迫使它改变这种状态为止。

\textbf{惯性}——质点不受力时保持静止或匀速直线运动状态的性质，其大小用质量量度。

\textbf{力}——使质点改变运动状态的原因。

质点处于静止或匀速直线运动时：
\[
\sum \vec{F}_i = 0 \quad\text{(静力学基本方程)}
\]

\subsection*{2. 牛顿第二定律}
\[
\sum \vec{F}_i = \frac{\mathrm{d}(m\vec{v})}{\mathrm{d}t} = \frac{\mathrm{d}m}{\mathrm{d}t}\vec{v} + m\frac{\mathrm{d}\vec{v}}{\mathrm{d}t}
\]

当质量不随时间变化时：
\[
\sum \vec{F}_i = m\frac{\mathrm{d}\vec{v}}{\mathrm{d}t} = m\vec{a}
\]

直角坐标系分量形式：
\[
\sum F_{ix} = m\frac{\mathrm{d}^{2}x}{\mathrm{d}t^{2}},\quad
\sum F_{iy} = m\frac{\mathrm{d}^{2}y}{\mathrm{d}t^{2}},\quad
\sum F_{iz} = m\frac{\mathrm{d}^{2}z}{\mathrm{d}t^{2}}
\]

自然坐标系分量形式：
\[
\sum F_n = ma_n = m\frac{v^{2}}{\rho},\qquad
\sum F_{\tau} = ma_{\tau} = m\frac{\mathrm{d}v}{\mathrm{d}t}
\]

\textbf{讨论：}
\begin{itemize}
  \item 第二定律只适用于质点
  \item 质量不能当常量的两种情况：变质量问题（火箭、雨滴）；高速运动（相对论效应，$v > 10^{6}\,\text{m/s}$）
\end{itemize}

\subsection*{3. 牛顿第三定律}
\[
\vec{F} = -\vec{F}'
\]
\begin{itemize}
  \item 成对性——物体之间的作用是相互的
  \item 同时性——相互作用相互依存，同生同灭
  \item 适用于接触力；非接触力（如万有引力）存在延迟效应
\end{itemize}

\section*{二、力学中常见的几种力}

\subsection*{1. 万有引力}
\[
F = G\frac{m_1 m_2}{r^{2}},\quad G = 6.67\times10^{-11}\,\text{m}^{3}\!\cdot\!\text{kg}^{-1}\!\cdot\!\text{s}^{-2}
\]
矢量形式：
\[
\vec{F}_{21} = -G\frac{m_1 m_2}{r^{2}}\,\vec{r}^{\,0}
\]

说明：
\begin{itemize}
  \item 引力质量与惯性质量对同一物体总是相等
  \item 定律只直接适用于两质点
  \item 质量均匀分布球体与质点可直接使用公式
  \item 一般物体与质点间的引力需用微积分计算
\end{itemize}

重力（地球表面附近）：
\[
P = G\frac{Mm}{R^{2}}(1 - 0.0035\cos^{2}\phi)
\]
其中 $\phi$ 为地理纬度。

\subsection*{2. 弹性力}
两宏观物体接触发生微小形变时产生。

胡克定律（弹簧）：
\[
\vec{f} = -kx\,\vec{i}
\]

绳子受拉伸时内部出现弹性张力。

\subsection*{3. 摩擦力}

\textbf{静摩擦力：}两物体保持相对静止但有相对运动趋势时产生。
\[
f_{\text{max}} = \mu_0 N
\]
$\mu_0$ 为最大静摩擦系数，$N$ 为正压力。

\textbf{滑动摩擦力：}
\[
f = \mu N
\]
$\mu$ 为滑动摩擦系数。

\section*{三、牛顿运动定律的应用}

\subsection*{动力学问题分类}
\begin{enumerate}
  \item 已知运动情况 $\to$ 求力（微分问题）
  \item 已知力 $\to$ 求运动情况（积分问题）
  \item 已知部分运动与部分力 $\to$ 求未知部分
\end{enumerate}

\textbf{加速度}是连接运动学与动力学的桥梁。

\subsection*{求解一般步骤}
\begin{enumerate}
  \item 选取研究对象（隔离法）
  \item 分析受力和运动情况，画受力分析图
  \item 建立坐标系
  \item 列方程求解（牛二定律 + 约束方程）
  \item 讨论分析结果的物理意义
\end{enumerate}

\subsection*{变力问题的处理方法}

\textbf{(1) 力随时间变化：$F = f(t)$}
\[
m\frac{\mathrm{d}v_x}{\mathrm{d}t} = f(t),\quad
v_x = \frac{1}{m}\int f(t)\,\mathrm{d}t + C
\]
再积分得运动学方程。

\textbf{(2) 力随速度变化：$F = f(v)$}
\[
\frac{\mathrm{d}t}{m} = \frac{\mathrm{d}v}{f(v)},\quad
t - t_0 = m\int\frac{1}{f(v)}\,\mathrm{d}v
\]

\textbf{(3) 力随位移变化：$F = f(x)$}
\[
f(x) = mv\frac{\mathrm{d}v}{\mathrm{d}x},\quad
\int f(x)\,\mathrm{d}x = \frac12 m(v^{2} - v_0^{2})
\]

\subsection*{变质量问题的处理方法}
牛顿第二定律的普适形式：
\[
\vec{F}_{\text{合}} = \frac{\mathrm{d}(m\vec{v})}{\mathrm{d}t}
= \frac{\mathrm{d}m}{\mathrm{d}t}\vec{v} + m\frac{\mathrm{d}\vec{v}}{\mathrm{d}t}
\]

\section*{四、牛顿运动定律的适用范围}

\subsection*{惯性系}
牛顿运动定律适用的参照系称为惯性系。
\begin{itemize}
  \item 地面参考系通常视为惯性系
  \item 相对于惯性系作匀速直线运动的参照系都是惯性系
\end{itemize}

\subsection*{适用范围}
牛顿运动定律适用于\textbf{宏观物体的低速运动}。
\begin{itemize}
  \item 高速运动 $\to$ 相对论力学
  \item 微观粒子 $\to$ 量子力学
\end{itemize}

\subsection*{惯性力（非惯性系）}
在非惯性系 $S'$ 中引入虚拟力——惯性力：
\[
\vec{F}_0 = -m\vec{a}_e
\]
其中 $\vec{a}_e$ 为 $S'$ 系相对惯性系的平动加速度。

在非惯性系中，牛顿第二定律形式上成立：
\[
\vec{F} + \vec{F}_0 = m\vec{a}_r
\]

\textbf{注意：}惯性力是虚拟力，没有施力者，也没有反作用力，不满足牛顿第三定律。

\end{document}
